\subsection{Experimental is not the same as functional: the distance from readout to activity}\label{sec:dd-assay}
\textbf{Claim.} Evidence codes say an experiment was done, not how close its readout sits to the
gene product's own activity. Expression changes (IEP), distal phenotypes and fractionation proteomics
are filed alongside reconstituted enzyme assays, and they are the main source of over-annotation
among experimental rows.
\textbf{Corpus.} IEP is the weakest experimental code (\accIEP\% accepted, \rejIEP\% rejected);
high-throughput direct assays (HDA) are rejected at \rejHDA\%; mutant phenotypes (IMP, IGI) are
accepted at only \accIMP\% / \accIGI\%, most of the rest kept as non-core.
\textbf{Projects.} \proj{ASSAY\_TO\_FUNCTION.md}: catalogue of 60 readout classes mined from cached
papers; molecular readouts license an MF term 77\% of the time (567/738) and phenotypic hub
readouts 8\% (90/1,087). \proj{BEHAVIOR.md}: 86\% of adjudicated behaviour annotations (169/197)
downgraded; accepted ones are sensory channels and receptors acting in the circuit.
\proj{IEP.md}: survey of all 25,401 IEP rows in UniProt-GOA. \proj{OMICS\_EVIDENCE.md}: locations
read from organelle proteomes hold up (63--89\% accepted) but vesicle and body-fluid proteomes do not
(6\%), and functions read from location or co-purification are unreliable under any code.
\textbf{Principle.} Necessity is not participation (the AGT/angiotensin-maturation comparator check).
\textbf{Message.} Downstream users who trust ``experimental'' and distrust ``electronic'' are
filtering on the wrong axis.
